\documentclass[11pt,letterpaper]{article}
\usepackage[margin=0.85in]{geometry}
\usepackage{amsmath,amssymb,booktabs,array}
\usepackage[T1]{fontenc}
\usepackage{lmodern,microtype}
\usepackage[hidelinks]{hyperref}
\setlength{\parindent}{0pt}
\setlength{\parskip}{6pt}
\setlength{\abovedisplayskip}{7pt}
\setlength{\belowdisplayskip}{7pt}
\begin{document}
\begin{center}
{\Large\bfseries Housing units, household choices, and equilibrium prices}\\[4pt]
{\small Tommaso De Santo --- explanatory note, October 1, 2026}
\end{center}

\textbf{The question.} Why does the model offer houses of size $2,4,6$ rather
than $0.2,0.4,0.6$, or $20,40,60$? The numerical scale is a choice of units.
It should not affect consumption, saving, tenure, fertility, or physical price
responses. Our equations admit that equivalence, provided quantities, prices,
and utility coefficients are converted together. Changing the grid alone would
change the economic model.

\textbf{What one housing unit means.} The empirical quantity anchor is the
American Housing Survey count of rooms, rather than bedrooms. The model uses
that room scale for housing quantities and assumes services proportional to
quantity, with a separate owner service premium. A four-room house is $H=4$
in room units. Define $h=H/10$ if one housing unit instead represents ten rooms;
the same house is then $h=0.4$. Its physical size has not changed. The empirical
room target must also be divided by ten when expressed in these new units.

\textbf{Consumption and affordability.} Income and consumption remain in the
same monetary units. If rent per room is $r$, rent per ten-room unit is $10r$.
The household therefore pays exactly the same amount:
\[
 rH=(10r)(H/10).
\]
The same identity applies to purchase value, down payments, collateral, resale
proceeds, depreciation, and property taxes. All these costs use price times
housing quantity, with unchanged financing and tax rates. Consequently the
same consumption and saving choices remain feasible. As an illustration,
not a calibration estimate:
\begin{center}
\begin{tabular}{lrr}
\toprule
 & Rooms & Ten-room units\\
\midrule
Same dwelling's quantity & $4$ & $0.4$\\
Rent per housing unit & $10$ & $100$\\
Total rent paid & $40$ & $40$\\
Income available & $100$ & $100$\\
Consumption after rent and a child expense of $10$ & $50$ & $50$\\
\bottomrule
\end{tabular}
\end{center}

\textbf{Fertility and utility.} Preserving the budget is necessary but not
sufficient: the value of consumption and housing must retain its comparison
with the value of children. For a fixed consumption weight $\alpha$, write
material utility as
\[
 u^{\mathrm{mat}}(c,H)
 =\frac{[c^\alpha H^{1-\alpha}]^{1-\sigma}}{1-\sigma}.
\]
Here $c$ is consumption and $\sigma$ is curvature of utility over the consumption--
housing composite. One transparent conversion replaces $H$ by $10h$ in this
expression. Material utility then stays exactly the same, and child utility
needs no alteration.

Owners use housing directly; renters use a closed-form allocation in which
the same unit conversion enters through rent.
An equivalent conversion is therefore to keep its functional form in $h$ and
multiply \emph{all utility-valued coefficients} by the common factor
$K=(0.1)^{(1-\alpha)(1-\sigma)}$. With $\alpha=0.733$ and $\sigma=2$,
$K=1.849$. This includes child benefits, the fixed birth cost, bequest strength,
and taste-shock scales. Total values then satisfy $V'=KV$, while the ratios of
value differences to taste-shock scales remain unchanged. Birth-attempt and
tenure probabilities, and optimal consumption and saving, are consequently
unchanged. Re-estimation is not needed for this conversion.

\newpage
\textbf{Why the derivative can change without changing behavior.} A derivative
is measured per numerical housing unit. An upgrade from four to six rooms is
$4\to6$ in room units and $0.4\to0.6$ in ten-room units. The numerical derivative
and the numerical increment change together. The finite utility comparison
for the same physical upgrade remains unchanged. For example, holding family
state and tenure fixed and ignoring a housing floor, the consumption sacrifice
at material indifference for $4\to6$ is $13.730\%$ under the baseline weight
$0.733$. It is also $13.730\%$ for $0.4\to0.6$ after a consistent conversion.

\textbf{Housing supply and price curvature.} Supply is expressed in the same
housing units as demand. Let $Q$ denote total housing, $H_0$ its supply scale,
$r$ unit rent, $\bar r$ reference unit rent, and $\xi$ supply elasticity:
\[
 Q=H_0(r/\bar r)^\xi,
 \qquad r=\bar r(Q/H_0)^{1/\xi}.
\]
Moving to ten-room units divides $Q$ and $H_0$ by ten and multiplies $r$ and
$\bar r$ by ten. The equation remains satisfied at exactly the same physical
allocation. Rescaling $H_0$ is part of this conversion; it does not require
making $H_0$ an endogenous equilibrium variable.

Inverse supply is convex when $\xi<1$. Nevertheless its proportional price
response is constant:
\[
 \frac{d\log r}{d\log Q}=\frac1\xi.
\]
At the specified $\xi=0.630$, a $5\%$ quantity expansion along the supply curve
requires an $8.052\%$ rent increase, in either set of units. The numerical
slope changes with the coordinates; the implied change in rent per actual room
does not. This calculation concerns the supply curve, not the full equilibrium
response to a preference or fertility shock, which also depends on demand and
the population closure.

\textbf{Application to the two implemented specifications.} The adopted
September 28 reference changes the housing weight with children and compensates
material utility using a fixed reference rent. That reference rent must also
be multiplied by ten. The compensation then restores the same common factor
$K$ across family states; scaling the grid and child benefit alone is insufficient.
In the later floor experiment, the housing weight is common and the physical
floor must be divided by ten along with the grid. Its population closure sets
population equal to total supplied housing divided by demand per household.
Both quantities convert together, leaving population and birth renewal unchanged.
These are separate specifications, not interchangeable calibration results.

\textbf{What the audit establishes.} A read-only review traced the implemented
budget, material utility, child benefit, bequest, choice-probability, aggregation,
and supply equations. Scalar arithmetic on transcribed formulas reproduced
the conversion identities at illustrative points, including family compensation
and the physical floor; it did not evaluate the renter closed-form allocation
or compare model outputs. Historical parameter conversions remain unchecked.
Source hashes cover the floor experiment, not the adopted reference.
No model was imported, recalibrated, or solved.
A full numerical equivalence certificate remains outstanding because the code
also contains fixed feasibility thresholds and numerical regularizers.

\textbf{Conclusion.} Using rooms is a coherent normalization: the numerical
label $2$ rather than $0.2$ does not itself select a different region of economic
curvature. What must be retained is the meaning of the house and every trade-off
it enters. This result does not validate proportional services per room,
continuous fractional rental quantities, financial affordability, or the imposed
demand and supply elasticities. Those are substantive assumptions and empirical
checks; a change of units neither fixes nor invalidates them. The source map and
arithmetic record are in the accompanying housing-unit audit.

\newpage
\begingroup
\setlength{\parskip}{4pt}
\begin{center}
{\large\bfseries From model prices to dollars: a worked example}
\end{center}

\textbf{The two units.} Housing is measured in rooms; money is measured relative
to the model's baseline annual gross earnings scale, which is normalized to one.
Let $M$ be the number of 2007 dollars corresponding to one model monetary unit.
Let $p$ be the purchase price of one room and $r$ its rent for one four-year model
period. The model links them by $r=up$, where $u$ is the period user-cost
coefficient. Neither $r$ nor $u$ determines the earnings-to-dollar conversion.

\textbf{Where the dollar scale comes from.} The diagnostic uses PSID 2007 gross
labor earnings of the reference person and spouse combined, not total household
income. For reference persons aged 18--65 with nonnegative earnings and positive
survey weights, it takes weighted means within twelve four-year age groups and
then averages those twelve means equally, matching the model income profile's
age-cell normalization. The result is \$77,963.06 in 2022 purchasing power. Using
the retained CPI-U annual indices converts this to 2007 dollars:
\[
 M\simeq77{,}963.06\times\frac{207.342}{292.655}
   \simeq\$55{,}235.75\quad\text{per model monetary unit}.
\]
This conversion is obtained from earnings, independently of the housing prices
being checked; it was not chosen to make those prices match the AHS.

\textbf{Purchase prices and rental payments.} Purchase value is a stock, so
multiply $p$ by $M$. Rent is a payment over four years, so multiply $r$ by $M$
and divide by 48 months. Using the full saved reference values gives:
\begin{align*}
 \text{Purchase value per room}
   &=0.7898695\ldots\times M=\$43{,}629.03,\\
 \text{Monthly rent per room}
   &=0.1423972\ldots\times M/48=\$163.86.
\end{align*}
The same conversion applied to the floor experiment gives the last column:
\begin{center}
\begin{tabular}{lrrr}
\toprule
Cost per room & AHS 2007 & Sept.\ 28 reference & Oct.\ 1 floor\\
\midrule
Monthly rent (\$) & 168.22 & 163.86 & 149.20\\
Purchase value (\$) & 43,495.49 & 43,629.03 & 39,723.80\\
\bottomrule
\end{tabular}
\end{center}
The data column uses weighted total annual contract rent divided by weighted
renter rooms and by twelve; owner value uses weighted total reported dwelling
value divided by weighted owner rooms. These are stock ratios, not averages of
individual households' price-per-room ratios.

\textbf{What the reference rent $\bar r$ does.} The supply equation uses a
separate, fixed number $\bar r=0.16$: $S(r)=H_0(r/\bar r)^\xi$. Thus $H_0$ is
the quantity supplied \emph{at that reference rent}, not necessarily at the
equilibrium rent $r$. On the same dollar conversion, $\bar r$ corresponds to
\$184.12 per room per month; households need not pay that amount. It is not a
price target or the earnings numeraire. It can be absorbed into the coefficient:
\[
 S(r)=Ar^\xi,\qquad A=H_0/\bar r^{\xi}\simeq19.966
 \quad(H_0=6.294,\ \xi=0.630).
\]
This is the same supply curve. A third reference, $r_*$ (code:
\texttt{utility\_reference\_rent}), compensates family-dependent utility weights
in the adopted reference; it is inactive in the floor specification. It is
distinct from both actual rent $r$ and supply reference rent $\bar r$.

\textbf{Scope of the comparison.} PSID earnings mainly concern the preceding
tax year; AHS rents and self-reported values are not quality-adjusted. The floor
point (no unsecured borrowing, chain 7, case 0173) remains experimental. Its preferences,
entry/credit rules and calibration procedure differ; the price gap does not
isolate one change. This preliminary check does not establish equivalence of
the calibration procedures. The
\href{price_level_readout.md}{price-level readout} and
\href{price_level_comparison.json}{saved calculation} retain full inputs and
qualifications.
\endgroup
\end{document}
